\boxde
\BTTN
\begin{ex}%[2D1B4-1]
    Phương trình đường tiệm cận ngang của đồ thị hàm số $y=\dfrac{x-3}{x-1}$ là
    \choice
    {$y=5$}
    {$y=0$}
    {$x=1$}
    {\True $y=1$}
    \loigiai
    {
        Ta có $\lim\limits_{x \to \pm \infty}\dfrac{x-3}{x-1} = \lim\limits_{x \to \pm \infty}\dfrac{1-\dfrac{3}{x}}{1-\dfrac{1}{x}}=1$, nên đường thẳng $y=1$ là tiệm cận ngang của đồ thị hàm số đã cho.
    }
\end{ex}


\begin{ex}
    Đường thẳng nào dưới đây là tiệm cận đứng của đồ thị hàm số $y=\dfrac{2x}{x+2}$?
    \choice
    {$x=2$}
    {$x=0$}
    {\True $x=-2$}
    {$x=1$}
    \loigiai{
        Tập xác định $\mathscr{D}=\mathbb{R}\setminus\{-2\}$.
        \begin{itemize}
            \item $\lim\limits_{x\to -2^+}\dfrac{2x}{x+2}=-\infty$.
            \item $\lim\limits_{x\to -2^-}\dfrac{2x}{x+2}=+\infty$.
        \end{itemize}
        Vậy $x=2$ là đường tiệm cận đứng của đồ thị hàm số.
    }
\end{ex}

\begin{ex}%[2D1Y4-1]
    Cho hàm số $y=\dfrac{x+1}{2x-2}$. Khẳng định nào sau đây đúng?
    \choice
    {Đồ thị hàm số có tiệm cận đứng là $x=\dfrac{1}{2}$}
    {Đồ thị hàm số có tiệm cận ngang là $y=-\dfrac{1}{2}$}
    {\True Đồ thị hàm số có tiệm cận ngang là $y=\dfrac{1}{2}$}
    {Đồ thị hàm số có tiệm cận đứng là $x=2$}
    \loigiai{
        Đồ thị hàm số $y=\dfrac{x+1}{2x-2}$ có tiệm cận đứng $x=1$ và tiệm cận ngang $y=\dfrac{1}{2}$.
    }
\end{ex}


\begin{ex}%[2D1Y4-1]
    Cho hàm số $y=f(x)$ có $\lim\limits_{x\to -\infty} f(x)= -2$ và $\lim\limits_{x\to +\infty} f(x)= 2$. Khẳng định nào sau đây đúng?
    \choice
    {Đồ thị hàm số đã cho có đúng một tiệm cận ngang}
    {Đồ thị hàm số đã cho không có tiệm cận ngang}
    {Đồ thị hàm số đã cho có hai tiệm cận ngang là hai đường thẳng $x=-2$ và $x=2$}
    {\True Đồ thị hàm số đã cho có hai tiệm cận ngang là hai đường thẳng $y=-2$ và $y=2$}
    \loigiai{
        $\lim\limits_{x\to -\infty} f(x)= -2$ nên $y=-2$ là tiệm cận ngang.\\
        $\lim\limits_{x\to +\infty} f(x)= 2$ nên $y=2$ là tiệm cận ngang.
    }
\end{ex}

\begin{ex}%[2D1Y4-1]
    Cho hàm số $y= \dfrac{2017}{x-2}$ có đồ thị $(H)$. Số đường tiệm cận của $(H)$ là
    \choice
    {$0$}
    {\True $2$}
    {$3$}
    {$4$}
    \loigiai{
        Đồ thị $(H)$ có tiệm cận đứng là $x=2$, tiệm cận ngang là $y=0$.\\
        Vậy số đường tiệm cận của $(H)$ là $2$.
    }
\end{ex}

\begin{ex}
    Tìm số đường tiệm cận của đồ thị hàm số $ y = \dfrac{x^2 - 3x + 2}{x^2 - 4}. $
    \choice
    {$1$}
    {$ 0$}
    {\True $2$}
    {$3$}
    \loigiai
    {
        Tập xác định: $ \mathscr D = \mathbb{R} \backslash \{\pm2 \} $.\\
        Ta có $ \lim \limits_{x \to \pm  \infty} y = 1 \Rightarrow  $ đồ thị hàm số có 1 tiệm cận ngang là $ y = 1. $\\
        Ta lại có $\lim \limits_{x \to 2} y =  \lim \limits_{x \to 2} \dfrac{x-1}{x+2} = \dfrac{1}{4} $ và $\lim \limits_{x \to -2^+} y =  \lim \limits_{x \to -2^+} \dfrac{x-1}{x+2} = -\infty$ nên đồ thị hàm số có 1 tiệm cận đứng là $ x = -2. $\\
        Vậy đồ thị hàm số đã cho có 2 đường tiệm cận.
    }
\end{ex}


\begin{ex}%[Đề minh họa BGD 2018-2019]%[2D1B4-1]
    \immini[thm]{Cho hàm số $y=f(x)$ có bảng biến thiên như sau. Tổng số tiệm cận ngang và tiệm cận đứng của đồ thị hàm số đã cho là
        \haicot
        {$4$}
        {$1$}
        {\True $3$}
        {$2$}
    }{
        \begin{tikzpicture}
            \tikzset{double style/.append style = {draw=\tkzTabDefaultWritingColor,double=\tkzTabDefaultBackgroundColor,double distance=2pt}}
            \tkzTabInit[nocadre=false, lgt=1.2, espcl=2.5,deltacl=0.6
            ]{$x$ /0.6,$y'$ /0.6,$y$ /1.6}{$-\infty$,$1$,$+\infty$}
            \tkzTabLine{,+, d ,+,}
            \tkzTabVar{-/ $2$ / , +D-/ $+\infty$ / $3$ , +/ $5$ /}
    \end{tikzpicture}}
    \loigiai{
        Từ bảng biến thiên ta có
        \begin{itemize}
            \item $\lim\limits_{x \to -\infty} y =2$ suy ra $y=2$ là tiệm cận ngang.
            \item $\lim\limits_{x \to +\infty} y =5$ suy ra $y=5$ là tiệm cận ngang.
            \item $\lim\limits_{x \to 1^-} y = +\infty$ suy ra $x=1$ là tiệm cận đứng.
        \end{itemize}
        Vậy đồ thị hàm số tổng cộng có $3$ đường tiệm cận ngang và tiệm cận đứng.

    }
\end{ex}

\begin{ex}%[Đề tập huấn Sở Ninh Bình, 2019]%[Nguyễn Văn Hải, dự án(12EX-5-2019)]%[2D1B4-1]
    \immini[thm]{Cho hàm số $y=f(x)$ có bảng biến thiên như hình bên. Hỏi đồ thị hàm số $y=f(x)$ có tổng số bao nhiêu tiệm cận (tiệm cận đứng và tiệm cận ngang)?
        \haicot
        {$0$}
        {\True $2$}
        {$3$}
        {$1$}
    }{
        \begin{tikzpicture}
            \tikzset{double style/.append style = {draw=\tkzTabDefaultWritingColor,double=\tkzTabDefaultBackgroundColor,double distance=2pt}}
            \tkzTabInit[nocadre=false,lgt=1,espcl=2,deltacl=0.6]
            {$x$ /0.6,$y’$ /0.6,$y$ /2.2}
            {$-\infty$ , $1$ , $3$ , $+\infty$}
            \tkzTabLine{,+,d,+,0,-,}
            \tkzTabVar{-/$-1$ ,+D- / $+\infty$ /  $-\infty$,+/ $2$, -/$-\infty$}
    \end{tikzpicture}}
    \loigiai{
        Ta có $\lim\limits_{x\to -\infty}f(x)=-1$, $\lim\limits_{x\to +\infty}f(x)=-\infty$ nên $y=-1$ là tiệm cận ngang.\\
        Ta có $\lim\limits_{x\to 1^+}f(x)=-\infty$ nên $x=1$ là tiệm cận đứng.\\
        Vậy đồ thị hàm số có $2$ đường tiệm cận.
    }
\end{ex}


\begin{ex}%[GHK1, THCS - THPT Nguyễn Khuyến, HCM, 2019]%[Vinh Vo, 12Ex3-2019]%[2D1B4-1]
    \immini[thm]{Cho hàm số $ y = f(x) $ xác định trên $ (-2;0) \cup (0;+\infty) $ và có bảng biến thiên như hình vẽ. Số đường tiệm cận của đồ thị hàm số $ f(x) $ là
        \haicot
        {$ 4 $}
        {$ 2 $}
        {$ 1 $}
        {\True $ 3 $}
    }{
        \begin{tikzpicture}
            \tikzset{double style/.append style = {double distance=2pt}}
            \tkzTabInit[lgt=1.2,espcl=3.5,nocadre=false]
            {$x$ /0.6, $f’(x)$ /0.7,$f(x)$ /1.7}
            { $-2$ , $0$ , $+\infty$}
            \tkzTabLine{d,+,d,-, }
            \tkzTabVar{D-/ /$ -\infty $, +D+/$ +\infty $/ $ 1 $, -/ $ 0 $}
            \draw[pattern = north west lines] ($(N13)-(0.2ex,0)$) rectangle (T11);
    \end{tikzpicture}}
    \loigiai{
        Từ bảng biến thiên, ta thấy $ \heva{& \lim \limits_{ x \to -2^{+} } f(x) = - \infty \\ & \lim \limits_{x \to 0^{-} } f(x) = + \infty \\ & \lim \limits_{x \to + \infty} f(x) = 0  } $, suy ra đồ thị hàm số $ f(x) $ có $ 3 $ tiệm cận trong đó có $ 2 $ tiệm cận đứng và $ 1 $ tiệm cận ngang.
    }
\end{ex}

\begin{ex}%[Đề thi khảo sát chất lượng trường THCS-THPT Lômônôxốp, Hà Nội 2018 ,Nhật Thiện 12EX1-2019]%[2D1Y4-1]
    \immini[thm]{Cho hàm số $y=f(x)$ xác định trên $\mathbb{R}\backslash \left\{{0}\right\}$, liên tục trên mỗi khoảng xác định và có bảng biến thiên như hình bên. Hỏi đồ thị hàm số có bao nhiêu đường tiệm cận?
        \haicot
        {$1$}
        {\True $2$}
        {$3$}
        {$4$}}{\begin{tikzpicture}
            \tikzset{double style/.append style = {draw=\tkzTabDefaultWritingColor,double=\tkzTabDefaultBackgroundColor,double distance=2pt}}
            \tkzTabInit[nocadre=false,lgt=1,espcl=2.3,deltacl=0.6]{$x$ /0.6,$y'$ /0.6,$y$ /1.8}{$-\infty$,$0$,$1$,$+\infty$}
            \tkzTabLine{,+,d,-,0,+}
            \tkzTabVar{+/ $2$ ,-D-/ $-\infty$/$-\infty$, +/ $1$ ,-/ $-\infty$ /}
    \end{tikzpicture}}
    \loigiai{
        Dựa vào bảng biến thiên, ta có $$\lim\limits_{x\to -\infty}y=2;\qquad \lim\limits_{x\to 0^{\pm}}y=-\infty$$
        Vậy hàm số có một tiệm cận ngang $y=2$, một tiệm cận đứng $x=0$.
    }
\end{ex}

\begin{ex}
    Số tiệm cận đứng của đồ thị hàm số $y=\dfrac{\sqrt{x+9}-3}{x^2+x}$ là
    \choice
    {$3$}
    {$2$}
    {$0$}
    {\True $1$}
    \loigiai{
        Tập xác định $\mathscr{D}=[-9;+\infty)\setminus \{-1;0\}$. \\
        Ta có $\left\{\begin{aligned}
            &\lim\limits_{x\to -1^+} \dfrac{\sqrt{x+9}-3}{x^2+x}=+\infty \\
            &\lim\limits_{x\to -1^-} \dfrac{\sqrt{x+9}-3}{x^2+x}=-\infty
        \end{aligned}\right. \Rightarrow x=-1$ là tiệm cận đứng. \\
        Ngoài ra $\lim\limits_{x\to 0} \dfrac{\sqrt{x+9}-3}{x^2+x}=\dfrac{1}{6}$ nên $x=0$ không là tiệm cận.}
\end{ex}

\begin{ex}
    Phương trình đường tiệm cận xiên của đồ thị hàm số $y=\dfrac{2x^2-x+1}{x-1}$ là
    \choice
    {$y=x-1$}
    {\True $y=2x+1$}
    {$y=2x+3$}
    {$y=x+1$}
    \loigiai{
        Sau khi chia đa thức, ta viết lại hàm số $y=2x+1+\dfrac{2}{x-1}$.\\
        Do $\lim\limits_{x\to \pm \infty}\left[y-(2x+1)\right]=\lim\limits_{x\to \pm \infty}\dfrac{2}{x-1}=0$ nên $y=2x+1$ là đường tiệm cận xiên.}
\end{ex}

\begin{ex}
    Giao điểm của đường tiệm cận đứng và đường tiệm cận xiên của đồ thị hàm số $y=\dfrac{x^2-3x+5}{x-2}$ có tọa độ là
    \choice
    {$(2;3)$}
    {$(-2;1)$}
    {\True $(2;1)$}
    {$(-2;3)$}
    \loigiai{
        Sau khi chia đa thức, ta viết lại hàm số $y=x-1+\dfrac{3}{x-2}$.
        \begin{itemize}
            \item [$\bullet$] Đồ thị hàm số có tiệm cận đứng là $x=2$
            \item [$\bullet$] Do $\lim\limits_{x\to \pm \infty}\left[y-(x-1)\right]=\lim\limits_{x\to \pm \infty}\dfrac{3}{x-2}=0$ nên $y=x-1$ là đường tiệm cận xiên.
        \end{itemize}
        Giải hệ $\heva{&x=2\\&y=x-1} \Leftrightarrow \heva{&x=2\\&y=1}$. Suy ra, giao hai đường tiệm cận có tọa độ $(2;1)$.
    }
\end{ex}

\begin{ex}%[Đề thi giữa HK1, THPT Bình Sơn Đồng Nai, 2019]%[Phan Minh Tâm, dự án EX3]%[2D1B4-2]
    Tiệm cận đứng của đồ thị hàm số $ y=\dfrac{2x+1}{x-m} $ đi qua điểm $ M(2;5) $ khi $ m $ bằng bao nhiêu?
    \choice
    {$ m=-2 $}
    {$ m=-5 $}
    {$ m=5 $}
    {\True $ m=2 $}
    \loigiai{
        Với $m \ne -\dfrac{1}{2}$ đồ thị có tiệm cận đứng là đường thẳng $ x=m $. Tiệm cận đứng $ x=m $ đi qua $ M(2;5) $ khi chỉ khi $ m=2 $.
    }
\end{ex}

\begin{ex}%[GHK1, THPT Quế Võ 2-Bắc Ninh, 2019]%[TranTony,12EX2]%[2D1B4-2]
    Cho hàm số $ y = \dfrac{2x^2-3x+m}{x-m} $ có đồ thị $ (C) $. Tìm tất cả các giá trị của tham số $ m $ để $ (C) $ không có tiệm cận đứng.
    \choice
    {\True $ m = 0 $ hoặc $ m = 1 $}
    {$ m = 2 $}
    {$ m = 1 $}
    {$ m = 0 $}
    \loigiai{
        Đồ thị $ (C) $ không có tiệm cận đứng khi $ m $ là nghiệm của $ 2x^2-3x+m $
        \begin{align*}
            \Leftrightarrow 2m^2 - 3m + m = 0 \Leftrightarrow \hoac{& m = 0 \\& m = 1.}
        \end{align*}
    }
\end{ex}

\BTTF

\begin{ex}
    Cho hàm số $y=f(x)=\dfrac{3-2x}{x+1}$. Xét tính đúng sai của các khẳng định sau:
    \choiceTF
    {\True Tập xác định của hàm số là $\mathbb{R}\backslash\{-1\}$}
    {\True  Đồ thị hàm số có đường tiệm cận đứng là $x=-1$}
    {Đồ thị hàm số có đường tiệm cận ngang là $y=3$}
    {Hai đường tiệm cận (đứng và ngang) của đồ thị tạo với hai trục tọa độ một hình phẳng có diện tích bằng $3$}
    \loigiai{
        \begin{itemchoice}
            \itemch Điều kiện xác định $x+1 \ne 0 \Leftrightarrow x \ne -1$. Suy ra $D=\mathbb{R}\backslash\{-1\}$.
            \itemch Đồ thị hàm số có tiệm cận đứng là $x=-1$.
            \itemch Đồ thị hàm số có tiệm cận ngang là $y=\dfrac{-2}{1}=-2$.
            \itemch Hai đường tiệm cận (đứng và ngang) của đồ thị tạo với hai trục tọa độ một hình chữ nhật như hình vẽ
            \begin{center}
                \begin{tikzpicture}[smooth,samples=300,scale=0.8,>=stealth]
                    \draw[->] (-3,0)--(2,0) node[below]{$x$};
                    \draw[->] (0,-3)--(0,1) node[right]{$y$};
                    \draw (0,0) node[below right]{$O$};
                    \draw[pattern = north west lines] (0,0)--(0,-2)--(-1,-2)--(-1,0);
                    \draw (-3,-2)--(2,-2)node[above]{\scriptsize TCN $y=-2$} (-1,-3)--(-1,1)node[above]{\scriptsize TCĐ $x=-1$};
                    \draw[fill=black] (-1,0) circle(1.5pt) (-1,-2) circle(1pt) (0,-2) circle(1.5pt);
                    \node[right] at (0,-2.3) {$A$};
                    \node[left] at (-1,0.3) {$B$};
                \end{tikzpicture}
            \end{center}
            Diện tích hình chữ nhật này là
            $$S=OA \cdot OB=2 \cdot 1=2.$$
    \end{itemchoice}}
\end{ex}

\begin{ex}%[2D1K4]
    Cho hàm số $y=f(x)$ xác định trên $(-\infty;2) \backslash\{-2\}$ và có bảng biến thiên như hình vẽ dưới đây.
    \begin{center}
        \begin{tikzpicture}
            \tikzset{double style/.append style = {draw=\tkzTabDefaultWritingColor,double=\tkzTabDefaultBackgroundColor,double distance=2pt}}
            \tkzTabInit[nocadre=false,lgt=1,espcl=3]
            {$x$ /0.7,$y'$ /0.7,$y$ /2.1}
            {$-\infty$,$-2$,$0$,$2$,$+\infty$}
            \tkzTabLine{,+,d,-,0,+,d,}
            \tkzTabVar{-/$-2$/,+D+/ $3$ / $+\infty$,-/$-2$/,+D/ $+\infty$ / }
            \draw[pattern = north west lines] ($(N43)+(0.1ex,0)$) rectangle (T21);
        \end{tikzpicture}
    \end{center}
    Xét tính đúng sai của các khẳng định sau:
    \choiceTF
    {\True Hàm số có giá trị nhỏ nhất bằng $-2$}
    {Hàm số có giá trị lớn nhất bằng $3$}
    {\True Đồ thị hàm số có hai đường tiệm cận đứng là $x=-2$ và $x=2$}
    {Đồ thị hàm số có hai đường tiệm cận ngang là $y=-2$ và $y=3$}
    \loigiai
    {
        Căn cứ vào bảng biến thiên của hàm số, ta có
        \begin{itemchoice}
            \itemch Hàm số đạt giá trị nhỏ nhất bằng $-2$ khi $x=0$.
            \itemch Do $\lim\limits_{x\to 2^{-}}f(x)=+\infty$ nên hàm số không có giá trị lớn nhất.
            \itemch Do $\lim\limits_{x\to -2^{+}}f(x)=+\infty$ và $\lim\limits_{x\to 2^{-}}f(x)=+\infty$ nên đồ thị hàm số có hai đường tiệm cận đứng là $x=-2$ và $x=2$.
            \itemch Do $\lim\limits_{x\to -\infty}f(x)=-2$ nên đồ thị hàm số có hai đường tiệm cận ngang là $y=-2$.
    \end{itemchoice}}
\end{ex}

\begin{ex}
    Cho hàm số $y=f(x)=\dfrac{5 x^2+9 x+9}{x-4}$. Xét tính đúng sai của các khẳng định sau:
    \choiceTF
    {Tập xác định của hàm số là $\mathbb{R}\backslash\{-4\}$}
    {\True Đường tiệm cận đứng của đồ thị hàm số có phương trình là $x=4$}
    {\True Đường tiệm cận xiên của đồ thị hàm số có phương trình là $y=5x+29$}
    {Giao điểm hai đường tiệm cận của đồ thị hàm số có toạ độ là $(4;29)$}
    \loigiai{
        Hàm số được viết thành $y=5x+29+\dfrac{125}{x-4}$.
        \begin{itemchoice}
            \itemch Điều kiện $x-4 \ne 0 \Leftrightarrow x \ne 4$. Suy ra $D=\mathbb{R}\backslash\{4\}$.
            \itemch Đường tiệm cận đứng của đồ thị hàm số có phương trình là $x=4$
            \itemch Đường tiệm cận xiên của đồ thị hàm số có phương trình là $y=5x+29$
            \itemch Giải hệ $\heva{&x=4\\&y=5x+29}\Leftrightarrow \heva{&x=4\\&y=49}$. Giao điểm hai đường tiệm cận của đồ thị hàm số có toạ độ là $(4;49)$.
    \end{itemchoice}}
\end{ex}

\begin{ex}
    Cho hàm số $y=f(x)=\dfrac{x^2-4 x+7}{x-1}$. Xét tính đúng sai của các khẳng định sau:
    \choiceTF
    {\True Đường tiệm cận đứng của đồ thị hàm số có phương trình là $x=1$}
    {\True Đường tiệm cận xiên của đồ thị hàm số có phương trình là $y=x-3$}
    {\True Giao điểm hai đường tiệm cận của đồ thị hàm số có toạ độ là $(1 ;-2)$}
    {Diện tích tam giác tạo bởi đường tiệm cận xiên của đồ thị hàm số và hai trục toạ độ là $\dfrac{9}{4}$}
    \loigiai{
        Hàm số được viết thành $y=x-3+\dfrac{4}{x-1}$.
        \begin{itemchoice}
            \itemch Đường tiệm cận đứng của đồ thị hàm số có phương trình là $x=1$
            \itemch Đường tiệm cận xiên của đồ thị hàm số có phương trình là $y=x-3$
            \itemch Giải hệ $\heva{&x=1\\&y=x-3}\Leftrightarrow \heva{&x=1\\&y=-2}$. Giao điểm hai đường tiệm cận của đồ thị hàm số có toạ độ là $(1;-2)$.
            \itemch Giao của đường thẳng $d \colon y=x-3$ với các trục tọa độ lần lượt tại $A(0;-3)$ và $B(3;0)$.\\
            Diện tích tam giác $OAB$ là $S=\dfrac{1}{2} OA \cdot OB=\dfrac{9}{2}$.
    \end{itemchoice}}
\end{ex}

\BTTL

\begin{ex}%[2D1N4-1]Câu 1
    Cho hàm số $y=\dfrac{2x-1}{x+3}$. Gọi $x=m$ và $y=n$ lần lượt là đường tiệm cận đứng và tiệm cận ngang của đồ thị hàm số. Tính giá trị của biểu thức $P=\dfrac{2m-1}{n+3}$.\\
    \shortans[3]{$-1{,}4$}
    \loigiai{Ta có:\\
        $\bullet \underset{x \to -3}{\lim}\,y=\infty \Rightarrow x=-3$ là đường tiệm cận đứng.\\
        $\bullet \underset{x \to +\infty}{\lim}\,y=2$ và $\underset{x \to -\infty}{\lim}\,y=2 \Rightarrow y=2$ là đường tiệm cận ngang.\\
        Vậy $m=-3; n=2 \Rightarrow P=\dfrac{2\cdot(-3)-1}{2+3}=\dfrac{-7}{5}=-1{,}4$.}
\end{ex}

\begin{ex}%[2D1B4-3]
    Cho đồ thị $(C)\colon y=\dfrac{x-3}{x+2}$ có hai đường tiệm cận cắt nhau tại $I$. Với $O$ là gốc tọa độ, hãy tính độ dài đoạn thẳng $OI$ (làm tròn đến hàng phần trăm).\\
    \shortans[3]{$2{,}24$}
    \loigiai{
        Ta có tiệm cận đứng của đồ thị $(C)$ là $x=-2$ và tiệm cận ngang là $y=1$. Do đó $I(-2;1)$ là giao điểm của hai đường tiệm cận của đồ thị $(C)$.\\
        Ta có $OI=\sqrt{ (-2-0)^2+(1-0)^2}=\sqrt{5} \approx 2{,}24$.
    }
\end{ex}

\begin{ex}
    Nếu trong một ngày, một xưởng sản xuất được $x$ kilôgam sản phẩm thì chi phí trung bình (tính bằng nghìn đồng) cho một sản phẩm được cho bởi công thức:
    $$
    y=C(x)=\dfrac{50x+2000}{x}
    $$
    Đồ thị hàm số $C(x)$ có một đường tiệm cận ngang (khi $x \to +\infty$) là $y=y_0$. Giá trị $y_0$ bằng bao nhiêu?\\
    \shortans[3]{$50$}
    \loigiai{
        Ta có $\lim\limits_{x \rightarrow+\infty} \dfrac{50x+2000}{x}=\lim\limits_{x \rightarrow+\infty} \left(50+\dfrac{2000}{x}\right)=50$.\\
        Vậy đường thẳng $y=50$ là tiệm cận ngang của đồ thị hàm số.
    }
\end{ex}

\begin{ex}%[2D1C4-3]Câu 6
    Cho hàm số $y=\dfrac{x-2}{x^2-3mx+m}$ tìm $m$ để đồ thị hàm số có đúng một tiệm cận đứng. Biết tổng các giá trị của tham số $m$ có dạng phân số $\dfrac{a}{b}$, tính tổng $S=a+b$.\\
    \shortans[3]{$101$}
    \loigiai{Dễ thấy tử thức có một nghiệm là $x=2$ do đó để đồ thị hàm số có đúng một tiệm cận đứng thì phương trình $x^2-3mx+m=0$ có nghiệm kép hoặc có hai nghiệm phân biệt trong đó có một nghiệm bằng $2$.\\
        $\Rightarrow \hoac{&\Delta=0\\&\heva{&\Delta>0\\&4-3\cdot2m+m=0}} \Leftrightarrow \hoac{&9m^2-4m=0\\&\heva{&9m^2-4m>0\\&4-3\cdot2m+m=0}} \Leftrightarrow \hoac{&m=0\\&m=\dfrac{4}{9}\\&\heva{&\hoac{&m<0\\&m>\dfrac{4}{9}}\\&m=\dfrac{4}{5}}} \Leftrightarrow \hoac{&m=0\\&m=\dfrac{4}{9}\\&m=\dfrac{4}{5}}$\\
        Vậy tổng các giá trị của tham số $m$ bằng $\dfrac{56}{45} \Rightarrow S=101$.
    }
\end{ex}

\begin{ex}%[2D1H4-1]Câu 4
    Cho hàm số $y=\dfrac{x^2+2x-3}{x-2}$, đồ thị hàm số có đường tiệm cận xiên có dạng $(C) \colon y=ax+b$. Tính giá trị của biểu thức $P=\dfrac{a}{b}$.\\
    \shortans[3]{$0{,}25$}
    \loigiai{Ta xét $y=\dfrac{x^2+2x-3}{x-2}=x+4+\dfrac{5}{x-2} \Rightarrow (C) \colon y=x+4$ là đường tiệm cận xiên của đồ thị hàm số. Vậy $P=\dfrac{1}{4}$.}
\end{ex}

\begin{ex}
    Gọi $d$ là đường tiệm cận xiên của đồ thị hàm số $y=mx+4-3m+\dfrac{3}{x+2}$, $m$ là tham số. Đường thẳng $d$ luôn qua điểm cố định $M$. Tính độ dài đoạn $OM$, với $O$ là gốc tọa độ.\\
    \shortans[3]{$5$}
    \loigiai{
        Đường tiệm cần xiên của đồ thị là $y=mx+4-3m \Leftrightarrow (x-3)m+4-y=0$.\\
        Ta có $(x-3)m+4-y=0,\,\forall m \Leftrightarrow \heva{&x-3=0\\&4-y=0} \Leftrightarrow \heva{&x=3\\&y=4}$.\\
        Đường thẳng này luôn qua điểm cố định $M(3;4)$. Khi đó $OM=\sqrt{3^2+4^2}=5$.
    }
\end{ex}